% DOCUMENT COOKBOOK / SHARED SETUP. INPUT THIS FILE; DO NOT COMPILE IT.
% The ordinary document remains flowing text. Every optional feature below
% stays inactive unless explicitly used in a document or recipe.

% FOUNDATION 01: PAGE GEOMETRY AND TYPOGRAPHY (edit existing values)
% Keep a normal, flowing research manuscript. Free positioning is opt-in.
\usepackage[T1]{fontenc}
\usepackage{lmodern,microtype}
\usepackage[letterpaper,margin=1in]{geometry} % change paper/margins here
\usepackage{setspace}
\setlength{\parindent}{1.3em}
\setlength{\parskip}{0pt}

% FOUNDATION 02: SCIENTIFIC CONTENT AND COLORS
% Named palette: customize once, and every command using that color changes.
% In the document body use \color{Ink} etc. For section styling use titlesec.
% Recipe 32: mathematical symbols. amssymb loads amsfonts (\mathbb, \mathfrak);
% amsmath supplies additional operators and delimiters; bm supplies bold math.
\usepackage{amsmath,amssymb,amsthm,bm}
% DEFAULT EQUATION NUMBERING: (1.1), (1.2), (2.1), ...
% article: reset by numbered section; report/book: reset by chapter.
% For continuous (1), (2), ... numbering in an article, comment out
% the following conditional. Starred headings do not reset the counter.
\ifcsname c@chapter\endcsname
  \numberwithin{equation}{chapter}
\else
  \numberwithin{equation}{section}
\fi
\usepackage{graphicx,caption,subcaption,wrapfig,float,placeins}
\usepackage{booktabs,tabularx,array,multicol}
\usepackage{xcolor,tikz}
\usetikzlibrary{calc,positioning}
\usepackage{titlesec}
\usepackage{fancyhdr}
\usepackage{listings}
\definecolor{CUGold}{HTML}{CFB87C}
\definecolor{Ink}{HTML}{193848}
\definecolor{Teal}{HTML}{237C89}
\definecolor{Muted}{HTML}{5E7179}
\definecolor{Soft}{HTML}{EFF5F6}
\graphicspath{{figures/}}
\captionsetup{font=small,labelfont=bf}
\lstset{basicstyle=\ttfamily\footnotesize,frame=single,breaklines=true,
  columns=fullflexible,keepspaces=true,showstringspaces=false,
  keywordstyle=\color{Teal},commentstyle=\color{Muted}}

% OPT-IN REVISION NOTES: recipe 30. The blank document shows NO notes unless
% you place \todo or \todo[inline] in its body. Source-only % TODO comments
% never appear in the PDF. For a clean submission change the package line to
% \usepackage[disable,textsize=footnotesize]{todonotes}
% Do not import todonotes a second time in the main .tex file.
\setlength{\marginparwidth}{1.90cm} % margin width fits the default 1-inch margin
\setlength{\marginparsep}{0.20cm}
\usepackage[textsize=footnotesize]{todonotes}

% FOUNDATION 03: CITATIONS AND BIBLIOGRAPHY (change OPTIONS here, not elsewhere)
% style=numeric: numbered references; style=authoryear: author-year citations.
% sorting=none: citation order; sorting=nyt: name/year/title order.
% In DOCUMENT BODY, use \nocite{*} before \printbibliography to include
% every item in references.bib, even sources not cited in the text.
% For journal articles only: \nocite{*} then \printbibliography[type=article].
% Do NOT load biblatex a second time or mix with natbib/BibTeX.
\usepackage[backend=biber,style=numeric,sorting=none]{biblatex}
\addbibresource{references.bib}
% Optional shortcuts, only if wanted:
% \newcommand{\PrintAllReferences}{\nocite{*}\printbibliography}
% \newcommand{\PrintArticlesOnly}{\nocite{*}\printbibliography[type=article]}
\usepackage{aliascnt}
\usepackage[colorlinks=true,linkcolor=Ink,citecolor=Teal,urlcolor=Teal]{hyperref}
\usepackage[noabbrev,nameinlink]{cleveref}

% MATHEMATICAL ENVIRONMENTS (recipe 18; advanced customization in 31).
% Alias counters let cleveref distinguish types while sharing numbering.
% All numbered environments share a counter, reset by numbered section.
% Optional name: \begin{definition}[Relative Entropy] ... \end{definition}
\theoremstyle{plain} % bold heading, italic body
\newtheorem{theorem}{Theorem}[section]
\newaliascnt{lemma}{theorem}
\newtheorem{lemma}[lemma]{Lemma}
\aliascntresetthe{lemma}
\crefname{lemma}{lemma}{lemmas}
\Crefname{lemma}{Lemma}{Lemmas}
\newaliascnt{proposition}{theorem}
\newtheorem{proposition}[proposition]{Proposition}
\aliascntresetthe{proposition}
\crefname{proposition}{proposition}{propositions}
\Crefname{proposition}{Proposition}{Propositions}
\newaliascnt{corollary}{theorem}
\newtheorem{corollary}[corollary]{Corollary}
\aliascntresetthe{corollary}
\crefname{corollary}{corollary}{corollaries}
\Crefname{corollary}{Corollary}{Corollaries}

\theoremstyle{definition} % bold heading, upright body
\newaliascnt{definition}{theorem}
\newtheorem{definition}[definition]{Definition}
\aliascntresetthe{definition}
\crefname{definition}{definition}{definitions}
\Crefname{definition}{Definition}{Definitions}
\newaliascnt{example}{theorem}
\newtheorem{example}[example]{Example}
\aliascntresetthe{example}
\crefname{example}{example}{examples}
\Crefname{example}{Example}{Examples}
\newaliascnt{assumption}{theorem}
\newtheorem{assumption}[assumption]{Assumption}
\aliascntresetthe{assumption}
\crefname{assumption}{assumption}{assumptions}
\Crefname{assumption}{Assumption}{Assumptions}

\theoremstyle{remark} % italic heading, upright body
\newaliascnt{remark}{theorem}
\newtheorem{remark}[remark]{Remark}
\aliascntresetthe{remark}
\crefname{remark}{remark}{remarks}
\Crefname{remark}{Remark}{Remarks}
% amsthm also provides the unnumbered proof environment automatically.
% Independent counter: replace the three definition alias/declaration lines
% (newaliascnt, newtheorem, aliascntresetthe) with the following line;
% keep the crefname/Crefname lines:
% \newtheorem{definition}{Definition}[section]
% Continuous numbering: remove [section] from the theorem declaration.
% Optional unnumbered variant (declare only if needed):
% \theoremstyle{definition}
% \newtheorem*{definitionstar}{Definition}
\theoremstyle{definition} % restore the previous default declaration style

% FOUNDATION: customize an OPTIONAL precision-placed title page.
% First fontsize argument is the font size (pt); second is baseline spacing.
\newcommand{\DocTitleFont}{\color{Ink}\fontsize{25}{31}\selectfont\bfseries}
\newcommand{\DocSubtitleFont}{\color{Muted}\fontsize{14}{20}\selectfont}
\newcommand{\DocAuthorFont}{\color{Ink}\fontsize{13}{18}\selectfont}
\newcommand{\DocAffiliationFont}{\color{Muted}\fontsize{11}{15}\selectfont}

% PAGE coordinates: X from page left, Y downward from page top (cm).
% Page objects are overlays: they occupy NO flow space; use on title pages.
% \DocPageTextAt{X}{Y}{WIDTH}{TEXT}
\newcommand{\DocPageTextAt}[4]{%
  \begin{tikzpicture}[remember picture,overlay]
    \node[anchor=north west,inner sep=0pt,outer sep=0pt]
      at ([xshift=#1cm,yshift=-#2cm]current page.north west) {%
        \begin{minipage}[t]{#3}\vspace{0pt}#4\par\end{minipage}};
  \end{tikzpicture}%
}
% \DocPageImageAt{X}{Y}{WIDTH}{FILE}
\newcommand{\DocPageImageAt}[4]{%
  \begin{tikzpicture}[remember picture,overlay]
    \node[anchor=north west,inner sep=0pt,outer sep=0pt]
      at ([xshift=#1cm,yshift=-#2cm]current page.north west)
      {\includegraphics[width=#3]{#4}};
  \end{tikzpicture}%
}

% CONTAINER-LOCAL coordinates: X right, Y down from the local TikZ origin.
% Place these ONLY within a tikzpicture (usually inside a figure/minipage).
% \DocImageAt{X}{Y}{WIDTH}{FILE}
\newcommand{\DocImageAt}[4]{%
  \node[anchor=north west,inner sep=0pt,outer sep=0pt] at (#1,-#2)
    {\includegraphics[width=#3]{#4}};%
}
% \DocTextAt{X}{Y}{WIDTH}{TEXT}
\newcommand{\DocTextAt}[4]{%
  \node[anchor=north west,inner sep=0pt,outer sep=0pt] at (#1,-#2) {%
    \begin{minipage}[t]{#3}\vspace{0pt}#4\par\end{minipage}};%
}
% \DocCanvas{WIDTH}{HEIGHT}{CONTENT}: explicit space, cm; no clipping.
\newcommand{\DocCanvas}[3]{%
  \begin{tikzpicture}[x=1cm,y=1cm]
    \useasboundingbox (0,0) rectangle (#1,-#2);
    #3
  \end{tikzpicture}%
}
% \DocAutoCanvas{CONTENT}: auto-sized bounding box; no reserved height.
\newcommand{\DocAutoCanvas}[1]{%
  \begin{tikzpicture}[x=1cm,y=1cm]
    \path (0,0);#1
  \end{tikzpicture}%
}

% OPT-IN REPEATED FOOTER. The blank document does NOT enable this.
% \DocSetFooter{LEFT}{CENTER}{RIGHT}{RULE THICKNESS}
\newcommand{\DocSetFooter}[4]{%
  \fancyhf{}\fancyfoot[L]{\footnotesize #1}
  \fancyfoot[C]{\footnotesize #2}\fancyfoot[R]{\footnotesize #3}
  \renewcommand{\headrulewidth}{0pt}%
  \renewcommand{\footrulewidth}{#4}%
  \pagestyle{fancy}%
}
